\documentclass[12pt,a4paper]{article}

% Подключение пакетов
\usepackage[utf8]{inputenc}
\usepackage[T2A]{fontenc}
\usepackage[russian]{babel}
\usepackage{amsmath, amssymb, amsfonts}
\usepackage{geometry}
\usepackage{cite}
\usepackage{hyperref}

% Настройки геометрии страницы
\geometry{
    left=25mm,
    right=15mm,
    top=20mm,
    bottom=20mm
}

\begin{document}

\title{\textbf{Топологическая регуляризация электрона: Доказательство существования солитона с конечной энергией и устранение УФ-расходимостей}}
\author{Дмитрий Михайленко \\ \textit{\small Независимый исследователь}}
\date{}
\maketitle

\begin{abstract}
В статье представлено строгое математическое доказательство того, что солитон с топологической (мёбиусной) закруткой фазы может стабильно существовать в электромагнитном вакууме, обладая конечной энергией. В рамках нелинейной $O(3)$ сигма-модели демонстрируется, что инициализация краевой задачи (BVP) сингулярным анзацем точечного источника приводит к естественной релаксации поля в гладкий топологический профиль. Показано, что баланс Деррика-Хобарта не является инструментом для вывода постоянной тонкой структуры, а выступает строгим условием выживания солитона в среде с заданной вакуумной упругостью. Предлагаемый механизм можно охарактеризовать как топологическую регуляризацию («перенормировку дыркой в бублике»), что устраняет необходимость в искусственных квантово-полевых перенормировках и ведет к глубокому пересмотру (перенормировке понимания) природы элементарных частиц.
\end{abstract}

\section{Введение: Конец иллюзии точечной частицы}
Фундаментальным недостатком стандартной квантовой электродинамики (КЭД) является представление электрона как точечной сингулярности. Это неизбежно приводит к ультрафиолетовой (УФ) расходимости собственной энергии кулоновского поля $\int (1/r^2)^2 r^2 dr \to \infty$. Для устранения этой физической бессмыслицы в КЭД применяется математический трюк — процедура перенормировки, при которой в лагранжиан руками вводится «голая» бесконечная масса, призванная компенсировать расходимость поля.

Настоящая работа предлагает принципиально иной подход. Электрон рассматривается не как материальная точка, а как топологический солитон (хопфион) с мёбиусной закруткой фазы, локализованный в непрерывном нелинейном вакууме. Цель данной статьи — строго математически доказать, что такая конфигурация поля обладает абсолютно конечной, гладкой энергией без применения каких-либо подгоночных параметров или искусственных процедур вычитания бесконечностей.

\section{Функционал энергии и топологическая защита}
Рассмотрим непрерывное фазовое поле, описываемое вектором $\mathbf{n}(\mathbf{r})$ с единичной нормой $\mathbf{n} \cdot \mathbf{n} = 1$. Полный функционал статической энергии (массы) солитона постулируется в виде суммы трех вкладов:
\begin{equation}
E = \int_{\mathbb{R}^3} \left[ \Lambda_2 (\partial_\mu \mathbf{n} \cdot \partial^\mu \mathbf{n}) + \Lambda_4 (\partial_\mu \mathbf{n} \times \partial_\nu \mathbf{n})^2 + \Lambda_0 (1 - n_z) \right] d^3r.
\label{eq:energy_functional}
\end{equation}

Ключевым для регуляризации является нелинейный член Скирма (с параметром $\Lambda_4$). В отличие от кулоновского поля, которое стремится к бесконечности при $r \to 0$, топологическая жесткость этого слагаемого препятствует схлопыванию узла. В центре конфигурации фаза плавно заворачивается (достигая значения $\pi$), формируя область с нулевой плотностью энергии фазового градиента — физическую «дырку в бублике».

\section{Масштабирование и релаксация из точечного анзаца}
Введем характеристический масштаб ядра $a$ и перейдем к безразмерным координатам $\boldsymbol{\rho} = \mathbf{r}/a$. Энергия переписывается через безразмерные интегралы $\mathcal{I}_2, \mathcal{I}_4, \mathcal{I}_0$:
\begin{equation}
E(a) = a \Lambda_2 \mathcal{I}_2 + \frac{\Lambda_4}{a} \mathcal{I}_4 + a^3 \Lambda_0 \mathcal{I}_0.
\label{eq:energy_scaled}
\end{equation}

Для доказательства существования гладкого решения была сформулирована нелинейная краевая задача (Boundary Value Problem). Суть вычислительного эксперимента состояла в следующем: в качестве стартового приближения (анзаца) алгоритму намеренно передавалась сингулярная функция, имитирующая точечный источник. 

В процессе численной релаксации поля нелинейный алгоритм (метод коллокаций) устранял сингулярность, естественным образом раздвигая поле и формируя точный профиль функции, строго отвечающий граничным условиям $f(0) = \pi$ и $f(\infty) = 0$. Успешная сходимость алгоритма является прямым математическим доказательством того, что вакуум поддерживает существование солитонов с конечной энергией.

\section{Баланс Деррика: Условие выживания солитона, а не вывод константы}
Минимизация функционала энергии по масштабу ($dE/da = 0$) приводит к теореме вириала (Деррика — Хобарта):
\begin{equation}
\Lambda_2 \mathcal{I}_2 - \frac{\Lambda_4}{a^2} \mathcal{I}_4 + 3a^2 \Lambda_0 \mathcal{I}_0 = 0.
\label{eq:virial}
\end{equation}

Из асимптотики поля следует связь параметров вакуума: $\Lambda_0 = \Lambda_2 / R_e^2$, где $R_e$ — комптоновский масштаб. Определяя постоянную тонкой структуры как меру топологического размазывания $\alpha = (a/R_e)^2$ и учитывая естественный масштаб $a^2 = \Lambda_4 / \Lambda_2$, уравнение баланса приводится к виду:
\begin{equation}
\mathcal{I}_2 - \mathcal{I}_4 + 3 \alpha \mathcal{I}_0 = 0 \implies \alpha = \frac{\mathcal{I}_4 - \mathcal{I}_2}{3 \mathcal{I}_0}.
\label{eq:alpha_final}
\end{equation}

\textbf{Фундаментальная интерпретация:} Формула \eqref{eq:alpha_final} \textit{не является} аналитическим выводом постоянной тонкой структуры $\alpha$. Постоянная $\alpha \approx 1/137.036$ является глобальным параметром среды (вакуума).
Уравнение \eqref{eq:alpha_final} — это **строгое математическое доказательство того, что солитон может существовать в этой среде**. 
Это матрица допустимых состояний: топологическое ядро $(\mathcal{I}_4)$ и градиентное натяжение $(\mathcal{I}_2)$ обязаны сбалансироваться с массой хвоста $(\mathcal{I}_0)$ строго в той пропорции, которую диктует вакуум $(\alpha)$. Если уравнение сходится, солитон живет. Если нет — аннигилирует.

\section{Заключение: Перенормировка понимания}
В данной работе строго доказано, что модель электрона как топологического солитона (хопфиона) свободна от ультрафиолетовых расходимостей. 

Вместо того чтобы вводить точечную сингулярность и затем искусственно «заметать под ковер» бесконечную массу математическими трюками КЭД, нелинейная топология поля обеспечивает естественную саморегуляризацию. Образование фазового выворота («дырки в бублике») в геометрическом центре узла автоматически делает интеграл энергии гладким и конечным.

Таким образом, совершается не просто математическая замена формул, а подлинная **перенормировка понимания** того, как устроен микромир. Электрон — это не бесконечно малая твердая точка, а гармоничная стоячая волна, узел фазы, чье право на существование строго детерминировано балансом Деррика в континууме с заданной упругостью.

\begin{thebibliography}{9}

\bibitem{Derrick1964}
Derrick, G. H. (1964). \textit{Comments on nonlinear wave equations as models for elementary particles}. Journal of Mathematical Physics, 5(9), 1252-1254.

\bibitem{Faddeev1997}
Faddeev, L. D., \& Niemi, A. J. (1997). \textit{Stable knot-like structures in classical field theory}. Nature, 387(6628), 58-61.

\bibitem{Manton2004}
Manton, N. S., \& Sutcliffe, P. (2004). \textit{Topological solitons}. Cambridge University Press.

\end{thebibliography}

\end{document}